\section{The register machine}
\label{sec:universality}

Both the strong construction of Section~\ref{sec:strong} and the weak variant of Section~\ref{sec:weak}
assemble the same register machine from the pieces of Section~\ref{sec:automaton}. We record the assembly and
the reduction here.

\subsection{Registers and operations}

A register is a chain of flip-flop bits, holding a number in binary as in Table~\ref{tab:reg}. The locomotive
reads and writes it by the three operations a counter machine needs. Increment ripples a carry through the
bits. Decrement clears the lowest set bit and borrows downward. The zero test is read off the same run, a
decrement that finds no set bit to clear returns by a separate track, and that return is how the machine learns
the register held zero. Each operation is a path through tracks and the three switches, and by
Lemmas~\ref{lem:direction} to~\ref{lem:switches} the locomotive follows it and leaves the bits in the right
state.

\subsection{The program}

A register machine is a finite list of instructions, each an increment with a jump or a decrement that branches
on the zero test. Each instruction is a small sub-circuit that drives the locomotive to its register, performs
the operation, reads the zero test, and routes the locomotive on to the next instruction. Loops and jumps are
track that returns. The program is a finite pattern of track and switches.

We verified the assembly at the level of the railway. Registers and a sequencer wired into a two-register
program, run by the single locomotive, compute, a multiplication program returns the product of its inputs for
several inputs, in agreement with a reference register-machine interpreter \cite{pollard2026vibe}. We also ran a
register physically as the cellular automaton itself, a binary counter driven by the single locomotive counts
correctly (Section~\ref{sec:verification}).

\subsection{The reduction}

\begin{proposition}[Assembly]
\label{prop:assembly}
The motion and switch rules of Section~\ref{sec:automaton} realize, on any regular tiling, a register with
correct increment, decrement, and zero test, and a sequencer that wires a finite program. With a register
machine of two counters, which is universal \cite{minsky1967}, the locomotive computes any partial recursive
function.
\end{proposition}

If the registers are supplied as infinite, ultimately-periodic track, the configuration is infinite and the
universality is weak, which is the variant of Section~\ref{sec:weak}. If instead the registers self-extend by
the build rule (B1), the configuration is finite and the universality is strong, which is
Theorem~\ref{thm:strong}. The two differ only by whether the registers are given or grown.
